\chapter{Практика}
\label{ch:chap2}

\section{Формирование несущих колебаний}
Итоговый QAM сигнал формируется из двух несущих колебаний: -sin и cos. Зададим несущие колебания на частоте 50 Гц

\begin{figure}[H]
    \centering
    \includegraphics{carrier.png}
    \caption{Несущие колебания}
\end{figure}

\section{Формирование символа}
Символ представляет собой пару чисел I и Q (можно записать как комплексное число), числа берутся из множества ${\pm 1, \pm 3}$ и управляют
параметрами сигнала (фазой и амплитудой). Возьмем I=3 и Q=1.

\section{Формирование радиосигнала}
Итоговый радиосигнал имеет вид $S_c(t) = I(t)cos(\omega_ct)-Q(t)sin(\omega_ct)$. В нашем случае он будет выглядеть так:

\begin{figure}[H]
    \centering
    \includegraphics{rf.png}
    \caption{Радиосигнал}
\end{figure}

Можно заметить, что радиосигнал имеет другую амплитуду и другую фазу.

Если посмотреть на спектр сигнала, то увидим спектральную линию на частоте 50 Гц и отзеркаленную линию на частоте -50Гц:

\begin{figure}[H]
    \centering
    \includegraphics{rf_spec.png}
    \caption{Спектр радиосигнала}
\end{figure}

\section{Прохождение сигнала через канал}

Для имитирования помех будем использовать нормальное распределение с $\mu = 0$ и $\sigma = 0.1$. Сгенерируем выборку из нормального
распределения с такими параметрами и добавим в наш сигнал:

\begin{figure}[H]
    \centering
    \includegraphics{rx_signal.png}
    \caption{Сигнал на приеме}
\end{figure}

Видим, что сигнал искаженный.

\section{Демодуляция QAM сигнала}

Демодуляция QAM сигнала происходит точно так же, как и модуляция: тоже необходимо домножить сигнал на cos и -sin несущей частоты.
При таком умножении мы получим не I и Q, а $\frac{I}{2}$ и  $\frac{Q}{2}$, т.е необходим усилитель после выхода умножителя (в программе просто домножаем на 2).

\begin{figure}[H]
    \centering
    \includegraphics{p.png}
    \caption{Сигналы на выходе перемножителей}
\end{figure}

Также появляются лишние спектральные линии на частотах  $-2f_c$ и $2f_c$, от которых необходимо избавиться путем фильтрации:

\begin{figure}[H]
    \centering
    \includegraphics{spec2.png}
    \caption{Спектры сигналов на выходе перемножителя}
\end{figure}


\section{Проектирование и созддание ФНЧ}

На рисунках выше видим лишние спектральные линии, нам нужно избавиться от них, взяв только центральную линию. Для этого нужен ФНЧ,
т.к он срезает верхние частоты и оставляет низкие. Выберем частоту среза равной 25 Гц, в нее как раз не попадут лишние частоты. \\

Далее сформируем импульсную характеристику ФНЧ (рисунок выше) и свернем ИХ с сигналами:

\begin{figure}[H]
    \centering
    \includegraphics{spec3.png}
    \caption{Спектры сигналов на выходе ФНЧ}
\end{figure}

\section{Анализ результата}
Спектр получился не очень аккуратный, но паразитные частоты исчезли. Теперь посмотрим, как наш сигнал изменился во временной области:

\begin{figure}[H]
    \centering
    \includegraphics{time.png}
    \caption{Итоговый демодулированный сигнал}
\end{figure}

Получили такой сигнал. Можем заметить, что синфазная часть сигнала колеблится примерно на значении 3, а квадратурная часть - примерно на значении 1.
Если усреднить сигнал, то получим как раз наши I=3 и Q=1. Далее можем подать I и Q на mapper и получить из них биты.

\endinput 